\chapter{Example Chapter}
\label{ch:example}

\begin{chapteroverview}
This chapter demonstrates a possible structure for notes and shows how
to use the formatting commands supplied by \texttt{formatting.sty}.
Delete or rename this file when beginning the real notes.
\end{chapteroverview}

\section{Learning objectives}

By the end of this chapter, the reader should be able to:

\begin{itemize}
    \item define the principal concepts;
    \item explain the relevant underlying mechanisms;
    \item recognise important clinical features;
    \item describe an appropriate approach to investigation; and
    \item summarise the broad principles of management.
\end{itemize}

\section{Definition}

\begin{definitionbox}[Example condition]
Insert a concise definition here. Define the condition before
discussing its causes, presentation, or management.
\end{definitionbox}

\section{Aetiology and risk factors}

Risk factors can be presented as a list:

\begin{itemize}
    \item risk factor one;
    \item risk factor two;
    \item risk factor three; and
    \item relevant protective factors.
\end{itemize}

\section{Pathophysiology}

Explain the underlying process in a logical sequence. Diagrams can be
stored in the \texttt{images} folder and inserted as follows:

\begin{verbatim}
\begin{figure}[htbp]
    \centering
    \includegraphics[width=\textwidth]{example-image}
    \caption{Description of the image.}
    \label{fig:example}
\end{figure}
\end{verbatim}

\section{Clinical presentation}

\subsection{Symptoms}

\begin{itemize}
    \item typical symptom;
    \item associated symptom; and
    \item relevant negative symptoms.
\end{itemize}

\subsection{Signs}

\begin{itemize}
    \item general examination finding;
    \item system-specific finding; and
    \item signs of complications.
\end{itemize}

\begin{redflag}
Record any features that may indicate a serious or time-critical
presentation.
\end{redflag}

\section{Investigations}

Tables work well for comparing investigations.

\begin{table}[htbp]
    \centering
    \caption{Example investigation table}
    \label{tab:example-investigations}

    \begin{tabularx}{\textwidth}{@{} l Y @{}}
        \toprule
        \textbf{Investigation} & \textbf{Purpose or expected finding} \\
        \midrule
        Bedside test
            & Add its purpose and interpretation here. \\

        Blood test
            & Add the relevant result or pattern here. \\

        Imaging
            & State when it is used and what it may demonstrate. \\
        \bottomrule
    \end{tabularx}
\end{table}

\section{Management}

Management can be divided into immediate, definitive, and longer-term
elements.

\subsection{Immediate management}

\begin{enumerate}
    \item Assess clinical stability.
    \item Address urgent problems.
    \item Arrange appropriate investigations and escalation.
\end{enumerate}

\subsection{Long-term management}

Summarise the principles of treatment, monitoring, prevention, and
follow-up.

\begin{clinicalnote}
Clinical recommendations should be checked against current local and
national guidance before being used in practice.
\end{clinicalnote}

\section{Complications}

\begin{itemize}
    \item early complications;
    \item late complications;
    \item treatment-related complications; and
    \item important effects on quality of life.
\end{itemize}

\section{Summary}

\begin{keypoint}[Key revision points]
\begin{itemize}
    \item Add the most important definition.
    \item Add the central mechanism.
    \item Add the characteristic presentation.
    \item Add the first-line investigation.
    \item Add the principal management approach.
\end{itemize}
\end{keypoint}

\lastupdated{2 September 2026}
